\documentclass[11pt]{article}
\usepackage[margin=1in]{geometry}
\usepackage{amsmath,amssymb,amsthm}
\usepackage{booktabs}
\usepackage{hyperref}
\usepackage{enumitem}
\usepackage{siunitx}

\hypersetup{colorlinks=true,linkcolor=blue,urlcolor=blue,citecolor=blue}

\title{Hegemon Prover Architecture (Cryptographer Version):\\
Formal Constraints, Security Accounting, Operation Budgets, and Feasibility}
\author{Hegemon Engineering}
\date{2026-02-28}

\begin{document}
\maketitle

\begin{abstract}
This document reformulates Hegemon proving as a constrained cryptographic systems problem. The goal is not to "optimize code" in the abstract; it is to satisfy explicit inequalities for security, latency, queue stability, and hardware budget. We derive operation counts from the current implementation constants (Poseidon2 rounds, trace dimensions, FRI settings, recursion payload shape), compute cycle budgets on \texttt{hegemon-prover} (32 cores, 256 GB RAM), and prove where the current architecture is infeasible. We then derive a replacement architecture from those constraints.
\end{abstract}

\section{Objective and Hard Constraints}

\subsection{Objective}
Sustain target throughput $T_{\mathrm{target}}$ with full post-quantum privacy, no proof bypass, and deterministic consensus import.

\subsection{Non-negotiable constraints}
\begin{enumerate}[label=C\arabic*.,leftmargin=2em]
  \item \textbf{Full PQ security:} no reduced parameters, no non-cryptographic fast path.
  \item \textbf{Fail-closed validity:} invalid proof data must be rejected at import.
  \item \textbf{Bounded pending latency:} no indefinite queue growth in steady-state.
  \item \textbf{Hardware reality:} heavy compute must run on \texttt{hegemon-prover}, not the laptop.
  \item \textbf{Acceptance gate:} sustained $\ge 10$ TPS for 30 minutes in full-PQ mode.
\end{enumerate}

\section{Implementation Constants (From Code)}

This section uses concrete constants currently in the codebase:

\begin{itemize}[leftmargin=2em]
  \item Transaction AIR rows: $N = 8192 = 2^{13}$.
  \item Trace width: $W=121$ columns.
  \item Poseidon2 parameters: width $t=12$, external rounds $R_F=8$ (split 4/4), internal rounds $R_P=22$, total steps per permutation $=31$.
  \item Merkle depth in transaction circuit: $32$.
  \item Per-transaction cycle schedule: cycle length $32$, total cycles $256$ (padded), used cycles $143$.
  \item Production FRI profile: log blowup $=4$ (blowup $b=16$), queries $q=32$.
  \item AIR max constraint degree used by quotient chunking logic: $8$, giving log quotient chunks $=\lceil\log_2(8-1)\rceil=3$, i.e. $8$ chunks.
\end{itemize}

Derived domain size for production FRI:
\begin{equation}
M = bN = 16\cdot 8192 = 131072 = 2^{17}.
\end{equation}

\section{Cryptographic Security Accounting}

\subsection{Inner proof soundness decomposition}
Write inner proof failure probability as
\begin{equation}
\varepsilon_{\mathrm{inner}} \le \varepsilon_{\mathrm{AIR}} + \varepsilon_{\mathrm{FRI}} + \varepsilon_{\mathrm{FS}} + \varepsilon_{\mathrm{hash}}.
\end{equation}
For FRI-like testing on quotient chunks with blowup $b$ and $q$ queries, the standard term is of the form
\begin{equation}
\varepsilon_{\mathrm{FRI}} \approx \left(\frac{1}{b}\right)^q
\end{equation}
up to low-order protocol constants.

For production parameters $b=16$, $q=32$:
\begin{equation}
\left(\frac{1}{16}\right)^{32}=2^{-128}.
\end{equation}
So the FRI term is in the 128-bit range before adding transcript/hash terms.

\subsection{Block-level recursion composition}
If one block aggregates $n$ statements and verifies one recursive root proof:
\begin{equation}
\varepsilon_{\mathrm{block}} \le \varepsilon_{\mathrm{outer}} + n\,\varepsilon_{\mathrm{inner}} + \varepsilon_{\mathrm{binding}} + \varepsilon_{\mathrm{hash}}.
\end{equation}
For $n\le 256$ and $\varepsilon_{\mathrm{inner}}\approx 2^{-128}$, the union term is $\le 2^{-120}$.

\subsection{Binding hash and commitment security}
Hegemon binds statements through 48-byte commitments and Blake-domain-separated hashes. With 384-bit outputs, collision resistance in a quantum model is still in the practical $\approx 2^{128}$ class (birthday-style quantum collision bounds), which is aligned with the target soundness class.

\section{Operator Cost Functions (Formal)}

Let TPS be $T$, block interval $\tau_b$, and transactions per block $n=T\tau_b$.

\subsection{Wallet}
\begin{align}
C_{\mathrm{wallet,send}} &= C_{\mathrm{select}} + C_{\mathrm{witness}} + C_{\mathrm{pack}} + C_{\mathrm{submit}},\\
C_{\mathrm{wallet,scan/sec}} &\approx T\cdot O\cdot C_{\mathrm{decrypt}},
\end{align}
where $O$ is outputs/tx.

\subsection{Coordinator}
\begin{equation}
C_{\mathrm{coord}}(m)\approx O(m\log m)+O(m\,C_{\mathrm{conflict}})+O(m\,C_{\mathrm{locality}}),
\end{equation}
for mempool size $m$.

\subsection{Prover (dominant)}
\begin{equation}
C_{\mathrm{prover}}(n)=C_{\mathrm{pre}}(n)+C_{\mathrm{inner}}(n)+C_{\mathrm{agg}}(n).
\end{equation}
Current implementation path for aggregation is monolithic in $n$ (one outer circuit that verifies all $n$ inner proofs), not a stage-tree recursion.

\subsection{Queue stability}
\begin{equation}
\frac{dQ}{dt}=\lambda-\mu,
\end{equation}
with stable operation requiring $\mu\ge\lambda$ over sustained windows.

\section{Exact Poseidon2 Operation Accounting (From Code)}

\subsection{Per-round operation counts}
From \texttt{poseidon2.rs}:
\begin{itemize}[leftmargin=2em]
  \item S-box $x\mapsto x^7$: 4 multiplications.
  \item External round (width 12):
  \begin{align}
  M_{\mathrm{ext}} &= 12\cdot 4 = 48,\\
  A_{\mathrm{ext}} &= 12\;\text{(add constants)} + 57\;\text{(mds\_light)} = 69.
  \end{align}
  \item Internal round:
  \begin{align}
  M_{\mathrm{int}} &= 4\;\text{(sbox on lane 0)} + 12\;\text{(diag mult)} = 16,\\
  A_{\mathrm{int}} &= 1\;\text{(lane-0 constant)} + 12\;\text{(sum)} + 12\;\text{(add sum)} = 25.
  \end{align}
  \item Initial linear step (step 0): $M_{0}=0$, $A_{0}=57$.
\end{itemize}

With 8 external rounds and 22 internal rounds:
\begin{align}
M_{\mathrm{perm}} &= 8\cdot 48 + 22\cdot 16 = 736,\\
A_{\mathrm{perm}} &= 57 + 8\cdot 69 + 22\cdot 25 = 1159.
\end{align}
Total field ops per permutation:
\begin{equation}
\mathrm{Ops}_{\mathrm{perm}} = M_{\mathrm{perm}} + A_{\mathrm{perm}} = 1895.
\end{equation}

\subsection{Per-transaction hash-step operations in witness generation}
The transaction trace loop runs 256 cycles, each with one Poseidon2 permutation schedule.
\begin{align}
M_{\mathrm{tx,witness}} &= 256\cdot 736 = 188{,}416,\\
A_{\mathrm{tx,witness}} &= 256\cdot 1159 = 296{,}704,\\
\mathrm{Ops}_{\mathrm{tx,witness}} &= 485{,}120.
\end{align}

Important: this is only witness-side hash algebra, \emph{not} full STARK proving.

\section{STARK Prover Lower-Bound Accounting}

\subsection{FFT/LDE lower bound}
For production domain $M=2^{17}$ and trace width $W=121$:
\begin{equation}
B_{\mathrm{trace}} = W\,M\,\log_2 M = 121\cdot 131072\cdot 17 = 269{,}615{,}104
\end{equation}
butterfly units.

Using a conservative 3 field ops per butterfly:
\begin{equation}
\mathrm{Ops}_{\mathrm{trace,LDE}}\gtrsim 808{,}845{,}312.
\end{equation}

Quotient chunk columns are 8 (from log chunks = 3), adding
\begin{equation}
\mathrm{Ops}_{\mathrm{quot,LDE}}\gtrsim 53{,}477{,}376.
\end{equation}

This is a lower bound for one transform pass; practical pipelines do multiple passes.

\subsection{Merkle commitment lower bound}
At least two major oracle trees (trace + quotient) plus FRI commit-phase trees. With $M=131072$ and observed commit-phase length $c=13$ (from recursion spike output):
\begin{align}
N_{\mathrm{nodes}}^{\mathrm{trace+quot}} &\approx 2(M-1)=262{,}142,\\
N_{\mathrm{nodes}}^{\mathrm{fri-phase}} &\approx \sum_{i=1}^{13}(M/2^i-1)=131{,}043,\\
N_{\mathrm{nodes,total}} &\approx 393{,}185.
\end{align}

If each compression uses one Poseidon2-like permutation cost envelope:
\begin{equation}
\mathrm{Ops}_{\mathrm{Merkle}}\gtrsim 393{,}185\cdot 1895=745{,}085{,}575.
\end{equation}

\subsection{Combined prover arithmetic floor}
\begin{equation}
\mathrm{Ops}_{\mathrm{prove,lb}} \gtrsim 8.09\times 10^8 + 5.35\times 10^7 + 7.45\times 10^8
\approx 1.61\times 10^9
\end{equation}
field-level operations before many runtime constants (memory movement, orchestration, serialization, recursion circuit plumbing).

\section{Cycle Budget on \texttt{hegemon-prover}}

For 32 cores at 2.5--3.0 GHz effective:
\begin{equation}
\mathrm{cycles/sec} \in [8.0,9.6]\times 10^{10}.
\end{equation}

Budget per tx:
\begin{align}
\mathrm{budget}_{1\mathrm{TPS}} &\in [8.0,9.6]\times 10^{10}\;\mathrm{cycles/tx},\\
\mathrm{budget}_{10\mathrm{TPS}} &\in [8.0,9.6]\times 10^{9}\;\mathrm{cycles/tx}.
\end{align}

Measured hot aggregation singleton is $S\approx 5.15$ s/tx, giving:
\begin{equation}
\mathrm{cycles}_{\mathrm{hot}}\in [4.12,4.94]\times 10^{11}\;\mathrm{cycles/tx}.
\end{equation}
That is $\approx 5.15\times$ over 1 TPS budget and $\approx 51.5\times$ over 10 TPS budget.

Measured cold miss equivalent $\sim 465$ s/tx is $\sim 465\times$ over 1 TPS budget.

\section{Where the Orders of Magnitude Are Hiding}

\subsection{Fact 1: not in core note hash math}
Witness-level hash ops per tx are $\sim 4.85\times 10^5$, tiny relative to total observed cycles.

\subsection{Fact 2: in proving-system and architecture constants}
Dominant terms are:
\begin{enumerate}[leftmargin=2em]
  \item LDE/FFT + Merkle commitment memory traffic and nontrivial constants.
  \item Recursion circuit setup/build and common-data preprocessing.
  \item Shape churn and cache misses.
  \item Monolithic aggregation path keyed by exact \texttt{tx\_count}.
\end{enumerate}

\subsection{Current shape-key entropy problem}
Current cache key includes
\begin{equation}
(\texttt{tx\_count},\texttt{pub\_inputs\_len},\texttt{log\_blowup},\texttt{proof\_shape}),
\end{equation}
so variable block occupancy induces many distinct expensive keys.

With observed cold build events in the hundreds of seconds, even rare misses destroy queue stability.

\section{Architectural Diagnosis (Mathematical)}

\subsection{Critical observation from code path}
Current aggregation for $n>1$ constructs one outer recursion circuit with a loop over all $n$ inner proofs. Tree metadata (arity/levels) is carried in payload, but proving work is monolithic per batch, not levelized DAG recursion.

Hence wall-time behavior is closer to:
\begin{equation}
T_{\mathrm{agg,current}}(n)\approx T_{\mathrm{cache}}(n)+T_{\mathrm{pack}}(n)+T_{\mathrm{runner}}(n)+T_{\mathrm{outer\_prove}}(n),
\end{equation}
with $T_{\mathrm{cache}}(n)$ potentially catastrophic on misses.

\subsection{Why this jams even with large hardware}
If cold events leak into critical path, effective service rate drops below arrival rate:
\begin{equation}
\mu_{\mathrm{effective}} < \lambda \Rightarrow Q(t)\ \text{diverges}.
\end{equation}
This is a control-stability failure, not just "insufficient cores".

\section{Replacement Architecture Derived from Constraints}

\subsection{Design theorem}
To make recursion scale, the system must satisfy:
\begin{enumerate}[leftmargin=2em]
  \item \textbf{Shape stationarity:} bounded shape set independent of instantaneous $n$.
  \item \textbf{Critical-path isolation:} cold preprocessing excluded from liveness path.
  \item \textbf{Levelized parallelism:} recursion executed as stage DAG (leaf/merge/root), not monolithic batch.
\end{enumerate}

\subsection{Tree-recursive model}
Let leaf capacity be $m$ tx/proof and merge arity $k$.
\begin{align}
L_0 &= \left\lceil\frac{n}{m}\right\rceil,\\
L_i &= \left\lceil\frac{L_{i-1}}{k}\right\rceil,\\
D &= \min\{d\,|\,L_d=1\}=\left\lceil\log_k L_0\right\rceil.
\end{align}

Parallel wall-time model:
\begin{equation}
T_{\mathrm{block}}\approx T_{\mathrm{pre}} + \max_j T_{\mathrm{leaf},j} + \sum_{i=1}^{D} T_{\mathrm{merge},i} + T_{\mathrm{root}}.
\end{equation}

Total work remains roughly linear in $n$, but critical path becomes logarithmic in tree depth rather than linear in proof count.

\subsection{Shape-count reduction}
Monolithic exact-$n$ caching requires up to $O(n_{\max})$ heavy shapes.
Tree design with fixed $(m,k)$ uses $O(\log_k n_{\max})$ merge shapes plus one leaf shape, dramatically reducing cold-synthesis risk.

\section{Feasibility Inequalities and Go/No-Go}

\subsection{Block-time feasibility}
For target TPS $T$ and block interval $\tau_b=6$:
\begin{equation}
T_{\mathrm{critical/block}} \le 6\ \mathrm{s},\qquad n=6T.
\end{equation}

\subsection{1 TPS gate}
Require all simultaneously:
\begin{align}
S_{\mathrm{hot}} &\le 1.0\ \mathrm{s/tx},\\
p95(S_{\mathrm{tx}}) &\le 2.0\ \mathrm{s/tx},\\
p_{\mathrm{cold,critical}} &\approx 0.
\end{align}

\subsection{10 TPS gate}
Require all simultaneously:
\begin{align}
T_{\mathrm{critical/block}} &\le 6\ \mathrm{s}\ \text{for}\ n\approx 60,\\
\mu_{30\mathrm{min}} &\ge 10\ \mathrm{TPS},\\
\sup_t\,Q(t) &\text{ bounded with finite drain time under burst tests.}
\end{align}

\subsection{Current decision}
\textbf{No-go} under present measured behavior.

Reason: measured hot and cold service times violate feasibility inequalities by large factors.

\section{Literature Review (Inspiration With Transfer Rules)}

\subsection{Transparent/PQ proof foundations}
\begin{itemize}[leftmargin=2em]
  \item STARK foundational framework \cite{stark2018}.
  \item DEEP-FRI soundness improvements \cite{deepfri2019}.
  \item Fractal transparent recursion \cite{fractal2019}.
\end{itemize}

\subsection{Modern recursive/folding line}
\begin{itemize}[leftmargin=2em]
  \item HyperNova folding architecture \cite{hypernova2023}.
  \item Neo/SuperNeo direction (small-field, pay-per-bit, PQ folding) from available metadata \cite{neo2025,superneo2026}.
\end{itemize}

\subsection{Permutation and implementation constants}
\begin{itemize}[leftmargin=2em]
  \item Poseidon2 design and round/linear-layer optimization \cite{poseidon22023}.
  \item Plonky3 recursion implementation model \cite{plonky3recbook}.
  \item RISC Zero recursion program model \cite{risc0repo,risc0docs}.
\end{itemize}

\subsection{Transfer rules for Hegemon}
From this literature, the transferable rules are:
\begin{enumerate}[leftmargin=2em]
  \item amortize expensive setup aggressively;
  \item keep recursion shapes fixed and cacheable;
  \item parallelize by proof DAG levels;
  \item treat memory bandwidth and cache locality as first-order constraints;
  \item enforce explicit soundness budgets under composition, not informal claims.
\end{enumerate}

\section{Actionable, Math-Driven Implementation Targets}

\subsection{Target envelope for first success (1 TPS)}
For $\tau_b=6$ s and $n=6$:
\begin{align}
T_{\mathrm{decode}} &\le 0.4\ \mathrm{s},\\
T_{\mathrm{commit}} &\le 1.8\ \mathrm{s},\\
T_{\mathrm{rec,critical}} &\le 2.8\ \mathrm{s},\\
T_{\mathrm{import}} &\le 0.7\ \mathrm{s},\\
T_{\mathrm{slack}} &\le 0.3\ \mathrm{s}.
\end{align}

\subsection{Architecture requirements implied by the envelope}
\begin{enumerate}[leftmargin=2em]
  \item Remove cold recursion setup from liveness path entirely.
  \item Replace exact-$n$ monolith with fixed-shape tree DAG.
  \item Promote prover cache from thread-local ownership to process-global manager with persistent hydration.
  \item Prebuild required shapes before entering throughput mode.
  \item Reject incompatible peers early (genesis/protocol/aggregation format).
\end{enumerate}

\section*{References}
\begin{thebibliography}{99}

\bibitem{stark2018}
E. Ben-Sasson, I. Bentov, Y. Horesh, M. Riabzev.
\newblock \emph{Scalable, transparent, and post-quantum secure computational integrity}.
\newblock IACR ePrint 2018/046.
\newblock \url{https://eprint.iacr.org/2018/046}

\bibitem{deepfri2019}
E. Ben-Sasson, L. Goldberg, S. Kopparty, S. Saraf.
\newblock \emph{DEEP-FRI: Sampling outside the box improves soundness}.
\newblock arXiv:1903.12243.
\newblock \url{https://arxiv.org/abs/1903.12243}

\bibitem{fractal2019}
A. Chiesa, D. Ojha, N. Spooner.
\newblock \emph{Fractal: Post-Quantum and Transparent Recursive Proofs from Holography}.
\newblock IACR ePrint 2019/1076.
\newblock \url{https://eprint.iacr.org/2019/1076}

\bibitem{hypernova2023}
A. Kothapalli, S. Setty.
\newblock \emph{HyperNova: Recursive arguments for customizable constraint systems}.
\newblock IACR ePrint 2023/573.
\newblock \url{https://eprint.iacr.org/2023/573}

\bibitem{neo2025}
W. D. Nguyen, S. Setty.
\newblock \emph{Neo: Lattice-based folding scheme for CCS over small fields and pay-per-bit commitments}.
\newblock IACR ePrint 2025/294 (metadata index).
\newblock \url{https://dblp.org/rec/journals/iacr/NguyenS25}

\bibitem{superneo2026}
W. Nguyen, S. Setty.
\newblock \emph{Neo and SuperNeo: Post-quantum folding with pay-per-bit costs over small fields}.
\newblock IACR ePrint 2026/242 listing metadata.
\newblock \url{https://eprint-classic.github.io/2026.html}

\bibitem{poseidon22023}
L. Grassi, D. Khovratovich, M. Schofnegger.
\newblock \emph{Poseidon2: A Faster Version of the Poseidon Hash Function}.
\newblock IACR ePrint 2023/323.
\newblock \url{https://eprint.iacr.org/2023/323}

\bibitem{plonky3recbook}
Plonky3 Contributors.
\newblock \emph{The Plonky3 recursion book}.
\newblock \url{https://plonky3.github.io/Plonky3-recursion/introduction.html}

\bibitem{risc0repo}
RISC Zero.
\newblock \emph{risc0 repository}.
\newblock \url{https://github.com/risc0/risc0}

\bibitem{risc0docs}
RISC Zero.
\newblock \emph{risc0-zkvm crate documentation}.
\newblock \url{https://docs.rs/risc0-zkvm/}

\end{thebibliography}

\end{document}
